\chapter{Graph-Design Screening}
\label{ch:experiments}

This chapter begins by verifying that a shared GNN actor can match the tuned
MLP baseline, then reports the graph-design screens. Screens A and B are the
completed seed-0 campaigns on the corrected information boundary across all screens; Screen E is a seed-0
factorial on the corrected boundary. All ask whether a specific choice in the
actor of Chapter~\ref{ch:methods} changes survival on held-out chronics.

\begin{tcolorbox}[
  colframe=darkred,
  colback=red!4!white,
  title=Information-boundary status of the reported screens
]
The preliminary viability runs and Screens C--D predate the corrected
local-graph construction of
Section~\ref{subsec:graph-local-global}. Their fixed candidate graphs marked
all rows of a neighboring substation as visible, including unattached busbars
and, in disaggregated graphs, neighboring private assets. Those isolated rows
could enter graph readout even when no live boundary line connected them to the
agent. Screens C--D therefore characterize the \emph{legacy,
information-leaking observation} and must not be presented as strictly
comparable to checkpoints trained with the corrected boundary. Screens A, B,
and E use connected context and controlled-only structural relations, and are
the leakage-free results reported in this chapter. The corrected Screens A and
B supersede the earlier campaigns under those names.
\end{tcolorbox}

\begin{itemize}
  \item \textbf{Screen A} varies the depth and width of the actor encoder,
  which applies to every graph schema: a $3\times4$ grid, 12 seed-0 runs.
  \emph{Complete.}
  \item \textbf{Screen B} varies which nodes enter graph readout and whether
  they are reduced by a mean or a maximum: a $4\times2$ factorial, seven new
  seed-0 runs plus one reused Screen-A cell. \emph{Complete.}
  \item \textbf{Screen C} varies the three structural augmentations of the
  busbar graph -- same-substation edges, substation-summary nodes, and
  virtual-node readout (Sections~\ref{sec:substation-representation}
  and~\ref{sec:virtual-node-readout}): a $2^3$ factorial, 24 runs.
  \emph{Complete.}
  \item \textbf{Screen D} varies the direction of the generator and load
  attachment relations in the disaggregated graph
  (Section~\ref{sec:heterogeneous-graph-representation}): a $3\times3$
  factorial, 27 runs. \emph{Complete.}
  \item \textbf{Screen E} varies candidate-node pooling, the action descriptor,
  and the idle-action head on the leakage-free disaggregated line-node graph: a $2^3$
  factorial, eight runs at seed 0, repeated with corrected graph inputs.
  \emph{Complete.}
\end{itemize}

\noindent Screens A and B form one corrected-boundary sequence on the busbar
graph and share the nominal \texttt{mp2\_h128} GINE/mean reference. Screen E is
the other corrected-boundary screen, on the disaggregated line-node graph. Screens C and
D are kept as separate legacy evidence and are not compared numerically with
the corrected campaigns.

\section{Preliminary GNN Encoder Viability}
\label{subsec:gnn-viability}

Before screening individual graph-design choices, an exploratory study checks
that a graph actor can match the tuned MLP baseline. This is a prerequisite for
transfer: the agents must be able to share one size-independent encoder without
sacrificing survival.

\paragraph{Question.}
Can a shared graph encoder perform as well as the MLP? If each agent requires a
separate encoder, or if the GNN cannot train stably, its structural
transferability is lost.

\paragraph{Baseline and changes.}
The study compares several GINE variants. It varies encoder size (light or
heavy), actor sharing (shared or unshared encoder), critic architecture (MLP or
GNN), and the training optimizations established for the MLP baseline,
including optimized critic updates and reward normalization.

\paragraph{Protocol.}
All GINE runs use \texttt{gnn\_type = gine}, $15{,}000{,}000$ total
timesteps, evaluation every $80{,}000$ timesteps, $2{,}000$ rollout steps,
deterministic evaluation, and the same 80/20 chronic split used by the MLP
experiments. GINE is used because it incorporates edge features and is
expressive over graph structures~\cite{xu2018how}. The complete graph-actor
pipeline is described in Appendix~\ref{app:gine-architecture};
Appendix~\ref{app:experiment-configurations} gives the full configurations in
Tables~\ref{tab:gine-run-configurations} and~\ref{tab:heavy-gine-configurations}.

\paragraph{Evidence.}
Table~\ref{tab:gine-rerun-summary} and Figure~\ref{fig:rerun-gine-best00}
compare the main architectural improvements with the \texttt{MLP Baseline}.
Figure~\ref{fig:gine-best-search2} shows the broader search over isolated
changes.

\begin{table}[H]
  \centering
  \small
  \caption{GINE rerun survival summary.}
  \label{tab:gine-rerun-summary}
  \resizebox{\textwidth}{!}{%
  \begin{tabular}{llrr}
    \toprule
    Run & Intended change & Final survival (\%) & Peak survival (\%) \\
    \midrule
    \texttt{Heavy GINE Baseline} & Historical GINE baseline & 59.2 & 93.9 \\
    \texttt{Light GINE (Concat, MLP Critic)} & Light GINE with MLP critic & 78.9 & 90.4 \\
    \texttt{No Bias, Opt Critic} & Optimized critic and bias 0.0 & 82.5 & 85.7 \\
    \texttt{Optimized Light GINE} & Light optimized GINE with MLP critic & 95.7 & 97.4 \\
    \bottomrule
  \end{tabular}
  }
\end{table}

\begin{figure}[H]
  \centering
  \includegraphics[width=\textwidth]{figures/rerun_gine_best00_vs_best10_14.png}
  \caption{Seed-aggregated test episodic survival comparing the
  \texttt{MLP Baseline} with high-performing GINE architectural variants.}
  \label{fig:rerun-gine-best00}
\end{figure}

\begin{figure}[p]
  \centering
  \includegraphics[width=\textwidth,height=0.9\textheight,keepaspectratio]{figures/gine_best_search2_baseline_comparisons.png}
  \caption{GINE survival comparisons. Each subplot compares the
  \texttt{MLP Baseline} with one variant from
  \texttt{configs/gine\_s0\_s1\_s2}. Full hyperparameters and architectural
  details are given in Appendices~\ref{app:gine-architecture}
  and~\ref{app:experiment-configurations}.}
  \label{fig:gine-best-search2}
\end{figure}

\paragraph{Interpretation.}
The broad search gives three conclusions:
\begin{itemize}
  \item \textbf{Shared encoder viability.} The unshared-actor variant does not
  outperform the shared variants, so one encoder can serve all agents without
  becoming the performance bottleneck.
  \item \textbf{Critic architecture.} A GNN centralized critic is unstable and
  degrades performance. The strongest variants pair the shared GNN actor with
  the standard MLP critic.
  \item \textbf{Optimization dependence.} Optimized critic updates, reward
  normalization, and removal of the initial do-nothing bias are also important
  for stable GNN training.
\end{itemize}

\noindent The strongest result is the \texttt{Optimized Light GINE}, which
combines the shared light GINE actor, the MLP critic, and the baseline
optimizations to reach $95.7\%$ final survival. A shared GNN is therefore
competitive on \texttt{bus14}, justifying the representation screens that
follow.


\section{How the Screens Are Organized}
\label{sec:screening-how-to-read}

Each screen is presented as one block containing its design and its outcome
together:

\begin{enumerate}
  \item \textbf{Question.} Which design choice is uncertain?
  \item \textbf{Design.} Which factors are varied, at which levels, and what is
  held fixed.
  \item \textbf{Result.} The measured tables.
  \item \textbf{Reading.} What follows, and what does not.
\end{enumerate}

\section{Shared Protocol}
\label{sec:screening-common-protocol}

Table~\ref{tab:screening-common-settings} separates the common optimization
protocol from the boundary and seed settings that differ by campaign. Within a
screen, every setting not named as a screened factor is fixed, so each screen
is a comparison against its own baseline cell.

\begin{table}[H]
  \centering
  \small
  \caption{Common optimization settings and campaign-specific boundary and
  seed settings.}
  \label{tab:screening-common-settings}
  \setlength{\tabcolsep}{4pt}
  \begin{tabular}{L{0.24\textwidth}L{0.68\textwidth}}
    \toprule
    Component & Fixed setting \\
    \midrule
    Environment & \texttt{bus14}, difficulty 0, topology actions, decentralized agents, normalized observations and rewards, 80/20 train/test chronic split with split seed 0 \\
    Actor & One GNN encoder shared by all agents; LayerNorm; node pre-encoder; two 128-unit action-head layers \\
    Information boundary & Screens A, B, E: connected context and controlled-only structural relations (corrected); Screens C, D: legacy whole-substation one-hop context \\
    Critic & Centralized MLP with three 256-unit hidden layers \\
    Graph preprocessing & None: neither deterministic physical scaling nor running standardization \\
    Training budget & \(15{,}000{,}000\) environment steps, 72 parallel environments, 576 rollout steps \\
    MAPPO update & 10 epochs, 16 minibatches, actor and critic learning rates \(10^{-4}\) with linear annealing, \(\gamma=0.9\), \(\lambda=0.95\), clip \(0.2\), target KL \(0.02\), entropy coefficient \(0.01\) \\
    Exploration & No initial do-nothing prior and no action-0 logit bonus \\
    Sparse control & Disabled: no intervention gate, no intervention penalty, no evaluation-time rho heuristic \\
    Evaluation & Deterministic evaluation every 82{,}944 environment steps, 10 episodes on each split \\
    Seeds & Screens A, B, C, D, E: 0 \\
    \bottomrule
  \end{tabular}
\end{table}

\subsection{Endpoint statistics}
\label{subsec:screening-endpoint-statistics}

Survival and the train/test split are as defined in
Section~\ref{sec:results-evaluation-protocol}: survival is the percentage of an
episode's maximum length that the agent keeps the grid alive, reported on the
test split. What differs here is how a run is reduced to one number. Each run saves the checkpoint with the highest test survival seen during training. That checkpoint is then re-evaluated deterministically over all 201 held-out test chronics.
\\
\\
Every run reaches the 15M-step budget, to within a negligible margin, so the
screens do not report how far their runs got. All screens report one seed
per cell. Every seed-level
endpoint is evaluated over the same 201 held-out chronics.

\subsection{Compute}
\label{subsec:screening-budget-alignment}

The campaigns were executed on JED (CPU) and IZAR (GPU). Within each factorial,
the execution setup was held fixed across the compared configurations; the
three-seed legacy campaigns also distribute seeds across both systems.
\\
\\
Evaluation is deterministic at a fixed checkpoint and seed. The corrected
screens have only one training seed per cell, so their
differences are descriptive rather than estimates of between-seed effects.
\\
\\
What limits the campaign is simulation throughput rather than model size. A
training step is dominated by stepping 72 Grid2Op environments, so the actor
forward and backward passes are a small part of the wall clock: across the depth
and width grid of Section~\ref{sec:screen-depth-width} the actor varies by a
factor of three in parameters, 83.6k to 244.9k, without a comparable effect on
run time. Parameter count is therefore not the cost to minimize when choosing an
architecture here.

\section{Screen A: Encoder Depth and Width}
\label{sec:screen-depth-width}

\subsection*{Question}

This screen measures how message-passing depth $K$ and encoder width $d_h$
behave after the information-boundary correction. On the busbar graph, depth
controls reach: one layer crosses one transmission-line edge, while two layers
let each busbar combine information across paths of two physical edges. The
graph-level readout then pools all valid node embeddings, so $K$ need not equal
the diameter of the local graph. Width controls how much information the node
and pooled embeddings can retain.
\\
\\
The result applies directly only to the busbar schema; using the same pair for
graphs with explicit asset or line nodes remains an assumption.

\subsection*{Design}

A complete $3\times4$ grid at seed 0, 12 runs.
Table~\ref{tab:depth-width-design} lists the screened parameters and their
levels. The graph output dimension is matched to the hidden dimension, so a
cell varies one width rather than two. Every other setting is inherited from
the corrected-boundary reference configuration of
Section~\ref{sec:screening-common-protocol}.

\begin{table}[H]
  \centering
  \small
  \caption{Screen A. Screened parameters and the values tested. All other
  settings are held at Table~\ref{tab:screening-common-settings}.}
  \label{tab:depth-width-design}
  \setlength{\tabcolsep}{5pt}
  \begin{tabular}{L{0.30\textwidth}L{0.24\textwidth}L{0.30\textwidth}}
    \toprule
    Parameter & Setting & Values tested \\
    \midrule
    Message-passing layers $K$ & \texttt{gnn\_layers} & 1, 2, 3 \\
    Encoder width $d_h$ & \texttt{gnn\_hidden\_dim} & 16, 32, 64, 128 \\
    \bottomrule
  \end{tabular}
\end{table}

\subsection*{Result}
The complete endpoint and marginal tables, the training curves, and the
parameter-count diagnostic are reported in
Appendix~\ref{app:screen-a-additional-results}.

\begin{figure}[H]
  \centering
  \includegraphics[width=0.86\linewidth]{figures/nl_s3dw_depth_width_heatmap.png}
  \caption{Screen A. The same twelve cells as a heatmap of the full-test
  survival of each run's best checkpoint. All cells use corrected graph inputs
  and seed 0.}
  \label{fig:depth-width-heatmap}
\end{figure}

\subsection*{Reading}

\textbf{The endpoint leader is \texttt{mp3\_h32}.} Its best checkpoint reaches
97.80\% full-test survival, 3.08 points above \texttt{mp2\_h128}
at 94.72\%.
\\
\\
\textbf{Depth does not separate cleanly.} The row means are 83.80\%, 82.90\%,
and 84.74\% for $K=1,2,3$. Their ranges overlap heavily, and the best and
second-worst cells both occur at $K=3$. These four single-seed cells per row do
not establish a reliable depth ordering.
\\
\\
\textbf{Width is strongly non-monotonic.} Averaged across depths, $d_h=32$
leads at 93.97\%, followed by 128 at 91.39\% and 16 at 88.42\%; every $d_h=64$
cell is weak, yielding a 61.49\% marginal mean.
\\
\\
\textbf{The factorial does not yield a simple architecture rule.} The depth
marginals differ by less than two points, width is non-monotonic, and the
ranking changes across depth--width combinations. The grid is therefore better
used to select a robust operating point than to claim a monotonic depth or
capacity effect.
\\
\\
Each endpoint covers all 201 test chronics, but it still reflects a checkpoint
selected from roughly 180 ten-episode evaluations and only one training seed.
Table~\ref{tab:depth-width-convergence} therefore complements it with metrics
computed from the complete training curves.

\begin{table}[H]
  \centering
  \small
  \caption{Screen A. Convergence measured from the training curves rather than
  from a single checkpoint. Steps to 90\% is the first evaluation at which the
  smoothed test survival reaches 90\%; the curve mean is the time-average of
  the smoothed curve over training. The final-window mean averages the raw last
  20 periodic evaluations, spanning approximately the final 1.58M environment
  steps. This is the campaign's late-stability window: it reduces sensitivity
  to the last few evaluations while remaining local to the converged regime.
  Dashes mark architectures that never reached the threshold.}
  \label{tab:depth-width-convergence}
  \setlength{\tabcolsep}{5pt}
  \begin{tabular}{lrrr}
    \toprule
    Architecture & Steps to 90\% (M) & Curve mean (\%) & Final-window mean (\%) \\
    \midrule
    \texttt{mp2\_h128} & \textbf{6.22} & \textbf{71.40} & \textbf{94.04} \\
    \texttt{mp1\_h16} & 6.88 & 70.92 & 84.93 \\
    \texttt{mp3\_h32} & 7.88 & 69.63 & 91.28 \\
    \texttt{mp2\_h16} & 8.79 & 69.16 & 88.21 \\
    \texttt{mp1\_h128} & 11.03 & 68.50 & 91.11 \\
    \texttt{mp1\_h32} & 6.64 & 63.13 & 93.77 \\
    \texttt{mp3\_h128} & -- & 58.58 & 80.20 \\
    \texttt{mp1\_h64} & -- & 51.60 & 63.37 \\
    \texttt{mp3\_h64} & -- & 49.80 & 67.58 \\
    \texttt{mp3\_h16} & -- & 48.82 & 60.00 \\
    \texttt{mp2\_h32} & 13.69 & 46.51 & 91.08 \\
    \texttt{mp2\_h64} & -- & 45.00 & 53.41 \\
    \bottomrule
  \end{tabular}
\end{table}

The curve summaries separate a best-checkpoint endpoint from sustained late
training behavior. \texttt{mp2\_h128} leads all three curve criteria: it has the
earliest 90\% crossing, the highest time-averaged curve, and the highest
20-evaluation final-window mean. Over that same final window its standard
deviation is 5.29 points, compared with 10.45 for \texttt{mp3\_h32}. The latter
still has the best full-test checkpoint, but its 3.08-point endpoint advantage
comes from one seed and does not persist across the late trajectory.
\\
\\
\textbf{Why carry \texttt{mp2\_h128} forward.} The selection combines the
training evidence with an architectural prior rather than treating the noisy
factor marginals as a law. In the busbar schema, $K=2$ is the smallest depth
that represents interactions across two consecutive transmission-line edges;
the global readout already combines all node embeddings, while a third layer
adds another round of mixing. This does \emph{not} mean that every pair of
nodes is at most two hops apart. Conditional on this deliberately chosen depth,
$d_h=128$ is the strongest width in the $K=2$ row both at full test and over the
late window, and its additional parameter cost is negligible beside Grid2Op
simulation. We therefore select $(K,d_h)=(2,128)$ as the conservative default
for the subsequent screens, while retaining \texttt{mp3\_h32} as the seed-0
endpoint challenger that would need a multi-seed confirmation.

\section{Screen B: Graph Readout Scope and Reducer}
\label{sec:screen-encoder-pooling}

\subsection*{Question}

The corrected local graph may contain nodes that provide message-passing
context without being equally useful in the final graph summary. Should readout
pool every visible node, only controlled nodes, only energized nodes, or the
intersection of controlled and energized nodes? For each scope, should the
valid embeddings be averaged or reduced by a per-channel maximum?
\\
\\
The scope mask is applied only at graph readout. Context nodes remain in the
message-passing graph, so this screen changes which embeddings form the policy
summary without changing what information can propagate to a controlled node.

\subsection*{Design}

A $4\times2$ factorial crosses four readout scopes---all visible, controlled,
energized, and controlled $\cap$ energized---with mean and max reduction at
seed 0. All cells use GINE with $K=2$ and $d_h=128$ on the corrected busbar
graph. Seven cells are new; the all-visible/mean cell reuses
\texttt{mp2\_h128} from Screen A because it is exactly the same configuration.

\subsection*{Result}

\begin{table}[H]
  \centering
  \small
  \caption{Screen B. Full-test survival by readout scope and reducer at seed 0.
  The all-visible/mean cell is inherited from Screen A.}
  \label{tab:encoder-pooling-grid}
  \setlength{\tabcolsep}{6pt}
  \begin{tabular}{lcc}
    \toprule
    & \multicolumn{2}{c}{Reducer} \\
    \cmidrule(lr){2-3}
    Readout scope & Mean & Max \\
    \midrule
    All visible & 94.72\textsuperscript{\dag} & 96.31 \\
    Controlled & 94.60 & 94.79 \\
    Energized & 98.93 & \textbf{99.39} \\
    Controlled $\cap$ energized & 91.45 & 96.47 \\
    \bottomrule
  \end{tabular}
  \\[0.4em]
  {\small \textsuperscript{\dag}Screen-A \texttt{mp2\_h128} reference; not a
  new run.}
\end{table}

\begin{table}[H]
  \centering
  \small
  \caption{Screen B training-curve diagnostics at seed 0, ranked by curve
  mean. ``To 90\%'' is the first crossing of 90\% by the five-evaluation
  trailing mean; curve mean is its time-weighted mean over training. Final-20
  statistics use the raw final 20 periodic evaluations (approximately 1.58M
  environment steps).}
  \label{tab:encoder-pooling-training-curves}
  \setlength{\tabcolsep}{4.5pt}
  \begin{tabular}{lrrrr}
    \toprule
    Scope/reducer & To 90\% (M) & Curve mean (\%) & Final-20 mean (\%) & Final-20 std (pp) \\
    \midrule
    Energized/max & \textbf{3.40} & \textbf{79.26} & 97.79 & 4.12 \\
    Energized/mean & 8.71 & 73.59 & \textbf{98.01} & \textbf{2.90} \\
    Controlled $\cap$ energized/mean & 5.97 & 72.97 & 91.18 & 8.59 \\
    All visible/mean\textsuperscript{\dag} & 6.22 & 71.40 & 94.04 & 5.29 \\
    Controlled/mean & 6.47 & 70.94 & 91.21 & 6.08 \\
    Controlled $\cap$ energized/max & 8.88 & 68.65 & 86.31 & 8.80 \\
    All visible/max & 4.73 & 68.38 & 88.57 & 8.08 \\
    Controlled/max & 10.87 & 68.27 & 90.67 & 6.81 \\
    \bottomrule
  \end{tabular}
\end{table}

The corresponding training curves are reported in
Appendix~\ref{app:screen-b-additional-results}.

\subsection*{Reading}

\textbf{Energized-node readout leads with both reducers.} Its mean-pooling cell
reaches 98.93\% and its max-pooling cell reaches 99.39\%, only 0.46 points
apart. Relative to pooling all visible nodes, energized-only pooling improves
survival by 4.21 points under mean and 3.08 points under max.
\\
\\
\textbf{Max is directionally consistent across scopes.} For every matched
scope, max pooling exceeds mean pooling on the full test. The gain is 1.59
points for all visible nodes, 0.19 for controlled nodes, 0.46 for energized
nodes, and 5.01 for the controlled-and-energized intersection. The magnitude
therefore depends strongly on scope, but the direction does not.
\\
\\
\textbf{Why carry energized/max forward.} It combines the best scope with the
reducer that wins every within-scope full-test comparison. Table~\ref{tab:encoder-pooling-training-curves}
also identifies it as the strongest overall training curve: it is the earliest
to cross 90\% (3.40M steps) and has the highest time-weighted curve mean
(79.26\%). Energized/mean has the better final-20 mean and standard deviation,
so the late-window evidence alone does not select max, and this remains a
seed-0 screening choice rather than a general claim that max pooling must
always dominate. Taken together with the 99.39\% full-test endpoint,
\texttt{energized\_max} is the most consistently supported operating point and
is fixed for the corrected reruns of Screens C and D, preventing readout from
becoming an additional factor in those experiments.

\section{Screen C: Structural Augmentations of the Busbar Graph}
\label{sec:screen-bus-structure}

\subsection*{Question}

Three optional augmentations add computational structure to the busbar graph
without adding any physical measurement: same-substation edges ($e$), which let
the busbars of one substation exchange messages directly; substation-summary
nodes ($n$), which insert one learned hub per substation; and virtual-node
readout ($v$), which replaces mean pooling with a learned graph-summary node.
Each is motivated by the same intuition, shorten the path between related
busbars, so the question is whether any of them earns its cost, and whether
they compose.

\subsection*{Design}

A complete $2^3$ factorial at seed 0, 8 runs. The graph type is
busbar throughout and the baseline cell is the unaugmented
\texttt{e0n0v0}. It has the same $K=2$, width-128, GINE/\texttt{energized\_max} design as
the corrected reference from Screen B. Cells are named by their three switches, so
\texttt{e1n0v1} means same-substation edges on, summary nodes off, virtual-node
readout on.

\subsection*{Result}

The detailed result tables and training curves are reported in
Appendix~\ref{app:screen-c-additional-results}.

\begin{figure}[H]
  \centering
  \includegraphics[width=0.86\linewidth]{figures/gs_s2_structure_bar.png}
  \caption{Screen C. The factorial results as a bar plot, with the baseline \texttt{e0n0v0} highlighted.}
  \label{fig:structure-cube}
\end{figure}

\subsection*{Reading}

The same-substation edges provide a slight improvement. The \texttt{e1n0v0} cell
is the best at 97.19\%, slightly outperforming the unaugmented \texttt{e0n0v0} baseline 
which achieved 95.45\%\footnote{This \texttt{e0n0v0} baseline is mathematically identical to the best-performing Energized/Max cell from Screen B, which achieved 99.39\%. The discrepancy in their final survival rates is due to the different hardware used for training (Screen B on Izar GPUs, Screen C on Jed CPUs), which affects the non-deterministic environment stepping and hardware-specific operations, leading to divergent trajectories.}. However, adding summary nodes or virtual-node readout still heavily 
degrades performance. Figure~\ref{fig:structure-curves} shows that the cells separate
well before the training budget ends.
\\
\\
The hierarchy factors have mixed effects. Summary nodes alone (\texttt{e0n1v0}) slightly improve the baseline to 96.09\%, but the virtual-node readout (\texttt{v1}) is particularly destructive, with cells containing it dropping to between 18\% and 77\%. When summary nodes and edges are combined (\texttt{e1n1v0}), performance degrades to 83.02\%.
\\
\\
The likely explanation is redundancy: both mechanisms add a learned
aggregation node, and together they stack two hubs between the busbars and the
readout (Figure~\ref{fig:virtual-node-readout}). With only two message-passing
layers, most physical nodes cannot reach that readout; Screen A's $K=3$ row
directly tests this interpretation.
\\
\\
Same-substation edges are milder and more consistent, and in the absence of other 
augmentations they successfully boost survival by about 1.74 points (from 95.45\% to 97.19\%). 
When mixed with the destructive hierarchy factors, their effect becomes highly erratic 
and often negative.
\\
\\
The subsequent stages therefore keep the plain busbar graph with \texttt{energized\_max} readout, as the 
modest gain from edges does not justify the additional graph density.
Any future hierarchy should be tested with a deeper encoder.

\section{Screen D: Generator and Load Message Direction}
\label{sec:screen-message-direction}

\subsection*{Question}

In the disaggregated graph, each generator and load is a node attached to its
busbar. Section~\ref{sec:heterogeneous-graph-representation} leaves the
direction of that attachment open: the asset may send messages to the busbar,
receive them from it, or both. The two directions are not symmetric in effect,
because the \texttt{energized\_max} readout pools
asset nodes as well as busbars. Does the choice matter, and do the generator
and load relations behave independently?

\subsection*{Design}

A complete $3\times3$ factorial over the generator and load directions at seed 0: 9 runs.
The graph type is disaggregated throughout; the line and summary
directions stay \texttt{bidirectional} and are inert for this schema. The
encoder matches the reference configuration: GINE, two layers, 128
features, \texttt{energized\_max} readout, no augmentations, so the only screened factors are
the two attachment directions. The baseline cell is bidirectional on both
relations. Directions are abbreviated \texttt{bi} (bidirectional), \texttt{a2b}
(asset to busbar), and \texttt{b2a} (busbar to asset).

\subsection*{Result}

The training curves and complete endpoint ranking are reported in
Appendix~\ref{app:screen-d-additional-results}.

\begin{figure}[H]
  \centering
  \includegraphics[width=0.96\linewidth]{figures/gs_hmd_direction_seed_heatmap.png}
  \caption{Screen D. Full-test survival of the best checkpoint by generator
  direction (rows) and load direction (columns). The baseline cell is
  \texttt{bi}/\texttt{bi}.}
  \label{fig:direction-seed-heatmap}
\end{figure}

\subsection*{Reading}

The optimal configuration requires asymmetric message passing directions for generators and loads.
The best configuration is \texttt{a2b} (asset-to-busbar) for generators and \texttt{b2a} (busbar-to-asset)
for loads, achieving a full-test survival of 99.12\%. Replacing the load direction with bidirectional (\texttt{bi})
yields a close second at 98.65\%. Conversely, the configuration previously identified as optimal under the corrupted data,
\texttt{b2a}/\texttt{b2a}, now performs poorly at 76.32\%.
\\
\\
The failure cases are equally clear. Using \texttt{bi} for generators and \texttt{b2a} for loads causes performance
to collapse to 59.98\%. Similarly, when both relations are asset-to-busbar (\texttt{a2b}/\texttt{a2b}), survival
is merely 73.75\%. The panel in Figure~\ref{fig:direction-curves} shows that these underperforming models separate
from the top configurations early in training.
\\
\\
A plausible mechanism is that
the \texttt{energized\_max} pooling includes generator and load nodes. By sending generator information to the busbars (\texttt{a2b})
but receiving information from busbars to loads (\texttt{b2a}), the network correctly routes production information
into the core topology, while allowing load nodes to incorporate their context before readout. This asymmetric flow
prevents the dilution of the pooled vector that occurs when nodes are either entirely inert or incorrectly pass messages
that blur their distinct roles.


\section{Screen E: Candidate-Action Scoring on the Corrected Boundary}
\label{sec:screen-e-cas}

\paragraph{Question.}
The candidate-action head of Section~\ref{sec:candidate-action-scoring} scores
every action from the encoded nodes it modifies rather than from a fixed-width
layer over an action index. Three choices in that head are unresolved: how the
touched nodes are pooled, whether the static action descriptor is appended to
the score, and whether the idle action gets its own head. This screen varies all
three.

\paragraph{Design.}
A $2^3$ factorial on \texttt{bus14}, seed $0$, eight runs. Pooling is
\texttt{mean} or \texttt{typed\_mean}; the action descriptor is appended
(\texttt{f1}) or withheld (\texttt{f0}); the idle action is scored by a
dedicated head (\texttt{a0h1}) or by the shared scorer (\texttt{a0h0}). Every
run uses the disaggregated line-node graph without one-hop context, a two-layer GINE
encoder of width $128$, and the 15M-step budget of
Section~\ref{sec:screening-common-protocol}. All eight reached 15M steps.

\paragraph{Result.}
Table~\ref{tab:screen-e-fulltest} reports the full-test survival of each best
checkpoint over the 201 held-out chronics. The corresponding training curves
and endpoint ranking are reported in
Appendix~\ref{app:screen-e-additional-results}.

\begin{table}[H]
  \centering
  \small
  \caption{Screen E, full-test survival of the best checkpoint on 201 held-out
  chronics. One seed per cell.}
  \label{tab:screen-e-fulltest}
  \setlength{\tabcolsep}{6pt}
  \begin{tabular}{lllrr}
    \toprule
    Pooling & Descriptor & Idle head & Survival (\%) & Best checkpoint \\
    \midrule
    \texttt{mean}       & \texttt{f0} & \texttt{a0h1} & 99.52 & 14{,}929{,}920 \\
    \texttt{typed\_mean} & \texttt{f1} & \texttt{a0h0} & 99.50 & 14{,}929{,}920 \\
    \texttt{mean}       & \texttt{f0} & \texttt{a0h0} & 99.18 & 14{,}929{,}920 \\
    \texttt{mean}       & \texttt{f1} & \texttt{a0h1} & 98.66 & 14{,}929{,}920 \\
    \texttt{mean}       & \texttt{f1} & \texttt{a0h0} & 98.32 & 14{,}846{,}976 \\
    \texttt{typed\_mean} & \texttt{f0} & \texttt{a0h1} & 98.17 & 14{,}764{,}032 \\
    \texttt{typed\_mean} & \texttt{f0} & \texttt{a0h0} & 92.50 & 12{,}358{,}656 \\
    \texttt{typed\_mean} & \texttt{f1} & \texttt{a0h1} & 57.84 & 11{,}612{,}160 \\
    \bottomrule
  \end{tabular}
\end{table}

\paragraph{Reading.}
\textbf{Six of the eight cells reach $98.2\%$ or better} and three exceed
$99\%$, so the candidate-action head is a viable actor on the corrected
boundary. No factor separates. The apparent marginals, \texttt{mean} at
$98.9\%$ against \texttt{typed\_mean} at $87.0\%$ and \texttt{f0} at $97.3\%$
against \texttt{f1} at $88.6\%$, are artifacts of two low cells; removing the
single failure at $57.84\%$ brings the pooling gap to about a point. At one seed
the defensible statement is that the head works and that the three choices do
not order themselves.
\\
\\
The descriptor is the one asymmetry worth recording: the best cell withholds it
and both failing cells include it, against a prior that favours the more
elaborate option. It is not redundant information, since every action still
reaches the scorer with its own pooled context -- agent~$0$'s $57$ actions
produce $57$ distinct row-sets, because an action toggles one to three of a
substation's lines rather than a single busbar. Withholding it removes a static
label and leaves the learned description intact.

\subsection{Replication with Corrected Graph Inputs}
\label{sec:screen-e-nls}

At seed 0, the eight cells were repeated with corrected graph inputs (NLS):
deterministic physical scaling and line-angle differences, without running
standardization. The GNN remains shared, but each agent retains its own scoring
MLP and optional idle head; this tests preprocessing, not scorer sharing.

\begin{table}[H]
  \centering
  \small
  \caption{Full-test survival over 201 held-out \texttt{bus14} chronics with
  original (NL) and corrected (NLS) graph inputs. The final column is NLS minus
  NL in percentage points. One seed per cell.}
  \label{tab:screen-e-nls-fulltest}
  \setlength{\tabcolsep}{4.5pt}
  \begin{tabular}{lllrrr}
    \toprule
    Pooling & Descriptor & Idle head & NL (\%) & NLS (\%) & $\Delta$ \\
    \midrule
    \texttt{mean}        & \texttt{f1} & \texttt{a0h0} & 98.32 & \textbf{100.00} & $+1.68$ \\
    \texttt{typed\_mean} & \texttt{f0} & \texttt{a0h0} & 92.50 & 99.84 & $+7.33$ \\
    \texttt{mean}        & \texttt{f0} & \texttt{a0h1} & 99.52 & 99.46 & $-0.07$ \\
    \texttt{mean}        & \texttt{f0} & \texttt{a0h0} & 99.18 & 99.43 & $+0.25$ \\
    \texttt{typed\_mean} & \texttt{f1} & \texttt{a0h1} & 57.84 & 90.71 & $+32.87$ \\
    \texttt{mean}        & \texttt{f1} & \texttt{a0h1} & 98.66 & 83.81 & $-14.85$ \\
    \texttt{typed\_mean} & \texttt{f1} & \texttt{a0h0} & 99.50 & 71.67 & $-27.82$ \\
    \texttt{typed\_mean} & \texttt{f0} & \texttt{a0h1} & 98.17 & 69.67 & $-28.50$ \\
    \bottomrule
  \end{tabular}
\end{table}

\begin{figure}[H]
  \centering
  \includegraphics[width=0.86\linewidth]{figures/candidate_action_preprocessing_fulltest.png}
  \caption{The paired comparison of
  Table~\ref{tab:screen-e-nls-fulltest}. NLS changes both power scaling and angle
  representation; the differences cannot be attributed to scaling alone.}
  \label{fig:candidate-action-preprocessing-fulltest}
\end{figure}

\paragraph{Reading.}
Four NLS cells exceed $98\%$, but three remain below $91\%$, with changes above
25 points in either direction. At one seed, this indicates an interaction with
the candidate head rather than a stable preprocessing effect. Only
\texttt{mean/f0} is robust: both idle-head choices exceed $99.1\%$ with NL and
$99.4\%$ with NLS. Chapter~\ref{ch:wcci-transfer} separately studies transfer
with a scorer shared across agents.
